%%
%% Copyright 2024 OXFORD UNIVERSITY PRESS
%%
%% v1.1 released August 2024
%%
%% This file is part of the 'pnas-nexus-authoring-template Bundle'.
%% ---------------------------------------------
%%
%% It may be distributed under the conditions of the LaTeX Project Public
%% License, either version 1.2 of this license or (at your option) any
%% later version.  The latest version of this license is in
%%    http://www.latex-project.org/lppl.txt
%% and version 1.2 or later is part of all distributions of LaTeX
%% version 1999/12/01 or later.
%%
%% The list of all files belonging to the 'pnas-nexus-authoring-template Bundle' is
%% given in the file `manifest.txt'.
%%
%% Template article for PNAS Nexus's document class `oup-authoring-template'
%% with bibliographic references
%%
%%
%% Change log
%%
%% v1.0 2021
%%    Adapted from oup_authoring_template.cls for use with pnas-nexus-authoring-template.tex
%% v1.1 August 2024
%%    Adjusted text of copyright statement for use with pnas-nexus-authoring-template.tex


\documentclass[unnumsec,webpdf,contemporary,large]{oup-authoring-template}%

%\usepackage{showframe}
\usepackage{booktabs}
\usepackage{lmodern}
\usepackage{anyfontsize}
\usepackage[dvipsnames]{xcolor}
\usepackage[table]{xcolor}

% line numbers
\usepackage[mathlines, switch]{lineno}
% \usepackage[right]{lineno}

\theoremstyle{thmstyleone}%
\newtheorem{theorem}{Theorem}%  meant for continuous numbers
%%\newtheorem{theorem}{Theorem}[section]% meant for sectionwise numbers
%% optional argument [theorem] produces theorem numbering sequence instead of independent numbers for Proposition
\newtheorem{proposition}[theorem]{Proposition}%
%%\newtheorem{proposition}{Proposition}% to get separate numbers for theorem and proposition etc.
\theoremstyle{thmstyletwo}%
\newtheorem{example}{Example}%
\newtheorem{remark}{Remark}%
\theoremstyle{thmstylethree}%
\newtheorem{definition}{Definition}


% \theoremstyle{thmstylethree}%
% \newtheorem{theorem}{Theorem}[section]
% \theoremstyle{definition}
% \newtheorem{definition}[theorem]{Definition}
% \theoremstyle{remark}
% \newtheorem{remark}[theorem]{Remark}
% \newtheorem{lemma}[theorem]{Lemma}
% \newtheorem{proposition}[theorem]{Proposition}
% \newtheorem{cor}[theorem]{Corollary}

\newcommand{\GF}{{\mathrm{GF}}}
\newcommand{\Z}{{\mathbb{Z}}}
\newcommand{\N}{{\mathbb{N}}}
\newcommand{\R}{{\mathbb{R}}}

\newcommand{\ACN}[1]{{\color{black} #1}}
\newcommand{\tolga}[1]{{\color{black} #1}}
\newcommand{\ergun}[1]{{\color{black} #1}}

\renewcommand\thefigure{S\arabic{figure}}
\renewcommand\thetable{S\arabic{table}}
%\usepackage{draftwatermark}
%\SetWatermarkText{Draft}

\title[Volumeteric Bravais  Weaves]{Supplementary Document: A Constructive  Framework for Discovery, Design, \& Classification of Volumetric Bravais Weaves}

\author[a]{Tolga T. Yildiz}
\author[b]{Alice C. Niemeyer}
\author[d,a]{Vinayak R. Krishnamurthy}
\author[c,a,1]{Ergun Akleman}

\authormark{Yildiz et al.}

\address[a]{\orgdiv{Department of Computer Science and Engineering}, \orgname{Texas A\&M University}, \orgaddress{\street{College Station}, \postcode{77840}, \state{TX}, \country{USA}}}
\address[b]{\orgdiv{Chair of Algebra and Representation Theory}, \orgname{RWTH Aachen University}, \orgaddress{\street{Pontdriesch 10-16}, \postcode{52062}, \state{Aachen}, \country{Germany}}}
\address[c]{\orgdiv{Visual Computing and Computational Media Section, School of Performance, Visualization and Fine Arts}, \orgname{Texas A\&M University}, \orgaddress{\street{College Station}, \postcode{77840}, \state{TX}, \country{USA}}}
\address[d]{\orgdiv{J. Mike Walker'66 Department of Mechanical Engineering}, \orgname{Texas A\&M University}, \orgaddress{\street{College Station}, \postcode{77840}, \state{TX}, \country{USA}}}


\begin{document}

\journaltitle{PNAS Nexus}
\DOI{DOI HERE}
\copyrightyear{YEAR}
\pubyear{YEAR}
\access{Advance Access Publication Date: Day Month Year}
\appnotes{Paper}

\firstpage{1}

%\subtitle{Subject Section}


% \leadauthor{Yildiz}


\corresp[$\ast$]{To whom correspondence should be addressed: \href{email:ergun.akleman@gmail.com}{ergun.akleman@gmail.com}}

\received{Date}{0}{Year}
\accepted{Date}{0}{Year}

%\editor{Associate Editor: Name}

\abstract{\ergun{This supplementary document provides additional details on the constructive framework for discovering, designing, and classifying volumetric Bravais weaves, specifically focusing on cP-weaves. We begin by establishing a mathematical representation of two-way two-fold weaves using elements of the finite field GF(2), demonstrating how these weaves can be systematically encoded within a lattice structure. The document further explores the transition from planar to volumetric weaves, detailing how unit cells can be represented using algebraic transformations and group-theoretic principles.
A computational methodology is introduced for enumerating regular plain cP-weaves, including a compact encoding scheme that leverages binary representations. Symmetry operations, including translation, rotation, and mirroring, are formally defined and applied to generate distinct weave configurations. The enumeration results, including the number of unique weave configurations under different symmetry constraints, are presented in tabular form for varying lattice sizes.
Additionally, a Python-based computational approach is provided, demonstrating efficient enumeration techniques using hash-table data structures to store and retrieve weave configurations. The results highlight the scalability of the framework and its potential for classifying and designing complex volumetric weaves.
By systematically formalizing volumetric weave structures, this document supports the main paper’s contributions to computational design, material science, and engineered structures, providing a foundation for future research in automated weave generation and classification.}
}
\keywords{
Volumetric Weaves, cP-Weaves, Woven Fabrics, Material Design
}


\maketitle


\firstpage{1}

%\subtitle{Subject Section}
% \leadauthor{Yildiz}
\onecolumn
\ergun{
\section{Representation of Two-way Two-fold Weaves as GF(2) Group on the Lattice Cells}

Symmetric two-way two-fold weaves have been formalized with crystallographic groups in \cite{periodic_geo_knots, O'Keeffe:cryst2-period} which subsumes the Bravais lattice in $\R^2$. Although two-way two-fold fabrics are studied thoroughly, the Bravais lattice based definition given in Definition~\ref{bravaisdefinition} does not have an inherent representation. In this section, we will give a methodology to represent the crossing information of local segments, i.e. the ranking order, with a cell-based representation. Doing so will enable us to derive a notation for two-way two-fold weaves and extend that notation to the volumetric weave of interest which are the cP-weaves.

We re-imagine the unit cells of a two-way two-fold fabric as a volumetric element that is derived from the Voronoi Cell of a tp lattice \cite{nebe2024catalogue} in $\mathbb{R}^2$ which results in a square with edges parallel to the standard basis on $\mathbb{R}^2$ (see Figure \ref{2dweaveunits}). 
%\ACN{To me this paragraph up to here is too repetitive and I would also like to move the sweeping curves to a later section. At this point we have the mathematical description and the symmetry  operations. How to translate this to a real world weave is not that obvious and deserves a longer explanation.} 
Note that two-way two-fold weaves can also be called tp-weaves to be consistent with cP-weaves. In $\mathbb{R}^2$ weave threads are connected from one parallel edge to the other parallel edge and then given an offset in the $z$ direction to give a rank order for the local configuration.

\begin{figure}[h]\setlength{\fboxsep}{0pt}
\centering
\captionsetup[subfloat]{width=0.99\linewidth}
\subfloat[A sample of a two-way two-fold plain weave.]{\includegraphics[width=0.35\linewidth]{Figures/1_1_1_top.pdf}}
\hspace{0.2in}
\captionsetup[subfloat]{width=0.99\linewidth}
\subfloat[Matrix encoding of plain weave.]{\includegraphics[width=0.35\linewidth]{Figures/1_1_1_matrix.pdf}}
\caption{Classical matrix view of 2-way 2-fold fabrics.\label{2dsubsume0}}
\end{figure}

\begin{figure}[h]\setlength{\fboxsep}{0pt}
\centering
\captionsetup[subfloat]{width=0.99\textwidth}%
\subfloat[Configurations of a single volume element for two-way two-fold weaves.]{
\label{2dweaveunits}
    \includegraphics[width=.25\textwidth]{Figures/0.pdf}  
    \includegraphics[width=.25\textwidth]{Figures/1.pdf}    
    }%
\hfill
\captionsetup[subfloat]{width=0.99\textwidth}%
\subfloat[Resulting representation of plain weave from the configurations.]{
    \label{subfig_2dweavedemo}
    \includegraphics[width=0.33\textwidth]{Figures/1_1_1.pdf}    
}\\
\caption{Demonstration of our framework for planar weaves. Here we can represent 2-way 2-fold fabrics as a result of a tp-lattice in 2D.\label{2dsubsume}}
\end{figure} 

Even though this appears to be a rather trivial exercise, it offers a direct pathway to cP-weave representation by providing us with a volumetric version of the matrix representation wherein each unit cell of the matrix is reinterpreted as a cube instead of a square (i.e. tp), and the ranking order is reinterpreted as a configuration of the threads, which can be encoded with matrix transformations. Note that, regardless of how we represent the threads inside a unit element (except that the threads lie in the same plane), there still exist only two distinct states (Figure ~\ref{2dweaveunits}) for two-way two-fold weaves. These two states can be transformed into one another using a transformation matrix that mirrors the assembly of the threads (about the $x-y$ plane in Figure ~\ref{2dweaveunits}). This gives us the cyclic group $GF(2)$ where each element is either the identity matrix $I$ or the mirror in $z$, $M_z$, as follows:
\begin{equation} 
\label{Eq:2DMatrices}
I =
\begin{pmatrix}
1 & 0 & 0\\
0 & 1 & 0\\
0 & 0 & 1\\
\end{pmatrix}
\; \; \; 
M_z =
\begin{pmatrix}
1 & 0 & 0\\
0 & 1 & 0\\
0 & 0 & -1\\
\end{pmatrix}
\end{equation}

The fundamental insight from the formalization using such a cyclic group is that it allows the representation of cP-weaves in terms of algebraic operations simply by replacing the group with one of its isomorphic groups. For instance, replacing the identity matrix $I$ with $0$, and the mirror matrix $M_z$ with $1$ turns the matrix multiplication into modulo $2$ addition since $I^2 = I$, $I M_z = M_z$, $M_z I = M_z$, and $M^2_z = I$. As a result these replacements directly lead us to the fundamental cyclic group $GF(2)$ that directly corresponds to the original cyclic group of two matrices. Furthermore, isomorphism offers a convenient way to encode volumetric weaving cases without the need to explicitly maintain ranking orders of the different threads.



To understand the algebraic operations through a concrete example, let us consider a periodic planar weave with a translation cell defined as an $N \times M$ grid made of squares in the plane. We can identify each square in this translation cell with two non-negative integers $(n,m)$ where $n=1,\ldots, N$; and $m=1,\ldots, M$ as $C_{(n,m)}$. Similar to matrix representation, assigning either $0$ or $1$ to each square for all $n$ and $m$ defines a weave. This can also be viewed as a discrete periodic function $V: \{1,\ldots,N\} \times \{1,\ldots,M\} \rightarrow \{0,1\}$ that can be represented as $V(n,m)$. We can further define discrete derivatives in the $x$ and $y$ directions as $\delta_x V = V(n+1,m)-V(n,m)$, and $\delta_y V = V(n,m+1)-V(n,m)$ respectively. Note that if $\delta_x=0$ or $\delta_y=0$ there is no change. But if $\delta_x=1$ or $\delta_y=1$ we take a mirror to get the orientation of the next cell, which corresponds to the rank order change in matrix representation. 

Here we can deduce the rank order for two-way two-fold weaves from the transformation matrix applied to the local cell. To see the effect of rank order change in a two-way two-fold weave, consider the threads shown in Figure~\ref{2dweaveunits}. The $x$-thread is the $y=0$ and $z=z_0$ constant line, and the $y$-thread is the $x=0$ and $z=-z_0$ constant line. If there is a change in the next unit cell along $x$, the $x$-thread becomes $y=0$ and $z=-z_0$ constant line, and the $y$-thread becomes $x=0$ and $z=z_0$ constant line where $z_0\neq 0$. Assume there is a change in $x$ direction, i.e. the change is between the cells $C_{(n,m)}$ and $C_{(n+1,m)}$. In such a case, the endpoints of two $x$-thread lines will not match in their $z$ coordinates. By connecting these two endpoints one can obtain a continuous curve. The same process is applied in the $y$-direction if there is a change. 
}

\section*{Computational Method for Enumeration of Regular Plain cP-Weaves}


To enumerate all the possible $N\times M\times L$ plain cP-weaves we need to use the compact representation:
\begin{equation}\label{eq:compactrep}
    \mathbf{W}=[x_0\dots x_{N-1}\;:\;y_0\dots y_{M-1}\;:\;z_0\dots z_{L-1}]
\end{equation}
Where each $x_i,y_j,z_k\in \GF(2), i\in [0,N-1],\;j\in[0,M-1],\;k\in[0,L-1]$. Which can be further simplified as 3 integers if we treat the collection of binary variables as a binary string that denotes an integer. This further simplifies the computational enumeration since all we need to define a plain cP-weave is three integers. 

Applying different isometries to the representation needs a definition on what are the operators that changes the representation. Given a $M\times M \times M$ regular plain cP-weave $\mathbf{W}_0=[x_0\dots x_{M-1}\;:\;y_0\dots y_{M-1}\;:\;z_0\dots z_{M-1}]$ we are going to define three symmetries. Note that $x_i'=1-x_i,y_j'=1-y_j, z_k'=1-z_k$.
\subsection*{Translation}
Translation cell of a plain cP-weave would cause negation of the entries that are not of the translation direction, i.e. an $x$-translation would cause negation of the representation of the $y$ and $z$ components and a circular shift of the $x$ component. All of these operations can be achieved with simple binary sequence operations. 
The translation by $x,\;y$ and $z$ would generate the following weaves:
\begin{align*}
    \mathbf{W}_0<<1_x=[x_1\dots x_{M-1}x_0\;:\;y_0'\dots y_{M-1}'\;:\;z_0'\dots z_{M-1}']\\
    \mathbf{W}_0<<1_y=[x_0'\dots x_{M-1}'\;:\;y_1\dots y_{M-1}y_0\;:\;z_0'\dots z_{M-1}']\\
    \mathbf{W}_0<<1_z=[x_0'\dots x_{M-1}'\;:\;y_0'\dots y_{M-1}'\;:\;z_1\dots z_{M-1}z_0]
\end{align*}
\subsection*{Rotation}
 Rotating the whole plain cP-weave would be another way to get a different geometry from the same weave topology. Particularly we can consider that all the rotations can be achieved by the combinations of rotations around the cardinal axes:
\begin{align*}
    \mathbf{W}_0\circ_{0,0,1:90}=[y_{M-1}'\dots y_0'\;:\;x_0'\dots x_{M-1}'\;:\;z_0\dots z_{M-1}]\\
    \mathbf{W}_0\circ_{0,1,0:90}=[z_{0}\dots z_{M-1}\;:\;y_{0}\dots y_{M-1}\;:\;x_{M-1}\dots x_{0}]\\
    \mathbf{W}_0\circ_{1,0,0:90}=[x_{0}\dots x_{M-1}\;:\;z_{M-1}'\dots z_{0}'\;:\;y_0'\dots y_{M-1}']\\
\end{align*}
\subsection*{Mirror}
Mirroring a plain cP-weave with respect to axis constant planes would result in the following representation change:
\begin{align*}
    \mathbf{W}_0|_x=[x_{M-1}'\dots x_0'\;:\;y_0'\dots y_{M-1}'\;:\;z_0'\dots z_{M-1}']\\
    \mathbf{W}_0|_y=[x_0'\dots x_{M-1}'\;:\;y_{M-1}'\dots y_0'\;:\;z_0'\dots z_{M-1}']\\
    \mathbf{W}_0|_z=[x_0'\dots x_{M-1}' \;:\;y_0'\dots y_{M-1}'\;:\;z_{M-1}'\dots z_0']
\end{align*}

\newpage
\section*{Python Code for Computational Method}

In the Python implementation we are not using the orbit finding algorithm since it would require more resources to compute. We employ a unordered hash-table data structure equivalent, i.e. \verb|set| data structure in Python, to store the weaves and use this data structure to find specific weaves in a give space of possible weaves. This data structure has $O(1)$ search which is going to be faster the lists of orbit finding algorithm. Another remark from the computational method is we can use three integers and their bits to represent a plain cP-weave, which reduces the memory costs as we are not using list of distinct bits.

\lstset{language=Python}
\begin{lstlisting}
# Utility functions for bit manipulations
# Reverses the bits of the N-bit integer
def reverseBits(n, N):
    return int('{:0{N}b}'.format(n, N=N)[::-1], 2)

# Rotates the weave w around x-axis by 90 degrees
def rotateX(w, N, M):
    return (w[0], M-reverseBits(w[2], N), M-w[1])

# Rotates the weave w around y-axis by 90 degrees
def rotateY(w, N, M):
    return (w[2], w[1], reverseBits(w[0], N))

# Rotates the weave w around z-axis by 90 degrees
def rotateZ(w, N, M):
    return (M-reverseBits(w[1], N), M-w[0], w[2])

# Mirrors the weave w about YZ plane
def mirrorX(w,N):
    return (2**N-1-reverseBits(w[0], N), 2**N-1-w[1], 2**N-1-w[2])

# Mirrors the weave w about XZ plane
def mirrorY(w,N):
    return ( 2**N-1-w[0], 2**N-1-reverseBits(w[1], N), 2**N-1-w[2])

# Mirrors the weave w about XY plane
def mirrorZ(w,N):
    return (2**N-1-w[0], 2**N-1-w[1], 2**N-1-reverseBits(w[2], N))

# Returns the derivative notation of weave w
def val_to_derivative(w, N):
    dx = ((w[0] << 1 | w[0] >> N-1)) ^ w[0]
    dy = ((w[1] << 1 | w[1] >> N-1)) ^ w[1]
    dz = ((w[2] << 1 | w[2] >> N-1)) ^ w[2]
    return [[w[0] >> N-1,w[1] >> N-1,w[2] >> N-1], [dx >> i & 1 for i in range(N-1,-1,-1)], [dy >> i & 1 for i in range(N-1,-1,-1)], [dz >> i & 1 for i in range(N-1,-1,-1)]]

# Calculates all the orbits for NxNxN regular plain cP Weaves
# rot, mirror and shift flags are used to control which symmetries are applied
def calculateTypes(N, rot=False, mirror=False, shift=False):
    c = 0
    num_sequences = 2**(N)
    M = num_sequences - 1
    type_classes = []
    # We are using sets to uniqely define weaves in a computation
    complete = set()
    # Adding all possible weave configurations to the set
    for i in range(num_sequences):
        for j in range(num_sequences):
            for k in range(num_sequences):
                complete.add((i,j,k))
    
    c = 0 
    # Rotation order for single pass rotation
    rotorder = ['x','x','z','x','x','z']

    # Going over the set until the set is empty
    while complete:
        c += 1

        # Get first weave w out of the set
        w = complete.pop()

        # Create a temporary set to generate the orbit of w
        tmp = set()
        if shift:
            i = w[0]
            j = w[1]
            k = w[2]

            # Shifting arrays to precompute shifts
            sx = list(map(lambda x: ((i << x) | (i >> N-x)) & (num_sequences-1), range(N)))
            nsx = list(map(lambda x: num_sequences-1-x,sx))
            sy = list(map(lambda x: ((j << x) | (j >> N-x)) & (num_sequences-1), range(N)))
            nsy = list(map(lambda x: num_sequences-1-x,sy))
            sz = list(map(lambda x: ((k << x) | (k >> N-x)) & (num_sequences-1), range(N)))
            nsz = list(map(lambda x: num_sequences-1-x,sz))
            for ii in range(N):
                for jj in range(N):
                    for kk in range(N):
                        x = sx[ii] if (jj + kk) % 2 == 0 else nsx[ii]
                        y = sy[jj] if (ii + kk) % 2 == 0 else nsy[jj]
                        z = sz[kk] if (jj + ii) % 2 == 0 else nsz[kk]
                        # w is shifted by ii, jj, kk in x, y, z direction respectively
                        w = (x,y,z)
                        tmp.add(w)

                        # Applying mirror, single mirror is enough to
                        if mirror:
                            tmp.add(mirrorZ(w,N))

                        # Applying rotations
                        if rot:
                            for r in rotorder:
                                for ir in range(4):
                                    w = rotateY(w, N, M)
                                    tmp.add(w)
                                    if mirror:
                                        tmp.add(mirrorZ(w,N))
                                if r == 'x':
                                    w = rotateX(w, N, M)
                                elif r == 'z':
                                    w = rotateZ(w, N, M)
        else:
            tmp.add(w)
            if mirror:
                tmp.add(mirrorZ(w,N))
            if rot:
                for r in rotorder:
                    for ir in range(4):
                        w = rotateY(w, N, M)
                        tmp.add(w)
                        if mirror:
                            tmp.add(mirrorZ(w,N))
                    if r == 'x':
                        w = rotateX(w, N, M)
                    elif r == 'z':
                        w = rotateZ(w, N, M)

        # Remove the orbit list from the all configurations
        complete -= tmp

        # Store the orbit in a list
        type_classes.append(tmp)

    return type_classes, c
\end{lstlisting}

\newpage
\section*{Tables for Enumerated Regular Plain cP-Weaves}

We have enumerated the possible configurations of regular plains with the computational method above. The Tables \ref{tab:222_1},\ref{tab:222_2},\ref{tab:222_3} enumerates the $2\times2\times2$ plain cP-weaves with different symmetries allowed. The Tables \ref{tab:444_1},\ref{tab:444_2} enumerates the $4\times4\times4$ plain cP-weaves with different symmetries allowed. We omitted the table enumerating the $4\times4\times4$ plain cP-weaves with only rotations because there are 240 orbits. The tables for plain cP-weaves after $4\times 4\times 4$ increases in size significantly, we decided not to include the tables but a report of the numbers in tables $\ref{tab:num1},\ref{tab:num2},\ref{tab:num3},\ref{tab:num4}$. This shows that the even though the plain cP-weaves is defined with constraints from the cP-weaves, it is able to express a rich space of shapes.


%%% Each figure should be on its own page

\clearpage

\begin{table}
    \centering
    \caption{Table of $2\times2\times2$ Regular Plain cP-Weaves with only rotational symmetries.}
    \begin{tabular}{ccccc}
    \midrule
        & Orbit Representative &  Visualization of $W$ & $|\textrm{Aut}(W)|$  & Orbit Length\\\toprule
         1 & $[10:10:11]$ &\includegraphics[width=1.2cm]{Figures/counting/222/223.png} & 2 & 12  \\\midrule
         2 & $[00:10:01]$ &\includegraphics[width=1.2cm]{Figures/counting/222/021.png} & 2 & 12  \\\midrule
         3 & $[11:11:01]$ &\includegraphics[width=1.2cm]{Figures/counting/222/331.png} & 2 & 12  \\\midrule
         4 & $[01:10:01]$ &\includegraphics[width=1.2cm]{Figures/counting/222/121.png} & 3 & 8  \\\midrule
         5 & $[00:11:01]$ &\includegraphics[width=1.2cm]{Figures/counting/222/031.png} & 4 & 6  \\\midrule
         6 & $[11:10:11]$ &\includegraphics[width=1.2cm]{Figures/counting/222/323.png} & 4 & 6  \\\midrule
         7 & $[11:11:00]$ &\includegraphics[width=1.2cm]{Figures/counting/222/330.png} & 8 & 3  \\\midrule
         8 & $[11:00:00]$ &\includegraphics[width=1.2cm]{Figures/counting/222/300.png} & 8 & 3  \\\midrule
         9 & $[00:11:00]$ &\includegraphics[width=1.2cm]{Figures/counting/222/030.png} & 24 & 1  \\\midrule
         10 & $[11:00:11]$ &\includegraphics[width=1.2cm]{Figures/counting/222/303.png} & 24 & 1  \\\midrule\midrule
          & & & Total & 64\\\bottomrule
    \end{tabular}
    \label{tab:222_1}
\end{table}
\clearpage
\begin{table}
    \centering
    \caption{Table of $2\times2\times2$ Regular Plain cP-Weaves with rotational and translational symmetries.}
    \begin{tabular}{ccccc}
    \midrule
         & Orbit Representative &  Visualization of $W$ & $|\textrm{Aut}(W)|$  & Orbit Length\\\toprule
         1 & $[10:00:10]$ &\includegraphics[width=1.2cm]{Figures/counting/222/202.png} & 8 & 24  \\\midrule
         2 & $[00:01:00]$ &\includegraphics[width=1.2cm]{Figures/counting/222/010.png} & 8 & 24  \\\midrule
         3 & $[01:10:01]$ &\includegraphics[width=1.2cm]{Figures/counting/222/121.png} & 24 & 8  \\\midrule
         4 & $[11:00:00]$ &\includegraphics[width=1.2cm]{Figures/counting/222/300.png} & 48 & 4  \\\midrule
         5 & $[11:11:00]$ &\includegraphics[width=1.2cm]{Figures/counting/222/330.png} & 48 & 4  \\\midrule\midrule
          & & & Total & 64\\\bottomrule
    \end{tabular}
    \label{tab:222_2}
\end{table}
\clearpage
\begin{table}
    \centering
    \caption{Table of $2\times2\times2$ Regular Plain cP-Weaves modulo rotational, mirror and translational symmetries.}
    \begin{tabular}{ccccc}
    \midrule
        & Orbit Representative &  Visualization of $W$ & $|\textrm{Aut}(W)|$  & Orbit Length\\\toprule
         1 & $[10:00:10]$ &\includegraphics[width=1.2cm]{Figures/counting/222/202.png} & 16 & 24  \\\midrule
         2 & $[00:01:00]$ &\includegraphics[width=1.2cm]{Figures/counting/222/010.png} & 16 & 24  \\\midrule
         3 & $[01:10:01]$ &\includegraphics[width=1.2cm]{Figures/counting/222/121.png} & 24 & 8  \\\midrule
         4 & $[11:11:00]$ &\includegraphics[width=1.2cm]{Figures/counting/222/330.png} & 24 & 8  \\\midrule\midrule
          & & & Total & 64\\\bottomrule
    \end{tabular}
    \label{tab:222_3}
\end{table}
\clearpage
\begin{table}
    \centering
    \caption{Table of $4\times4\times4$ Regular Plain cP-Weaves with rotational and translational symmetries. Red rows indicate that these cP-weaves can be obtained with repeating $2\times2\times2$ plain cP-weaves. Green rows indicate that these cP-weaves can be constructed from repeating $4\times4\times2$ plain cP-weaves. Blue rows indicate that these cP-weaves can be constructed from $4\times2\times2$ plain cP-weaves.}
    \begin{tabular}{ccccc}
    \toprule
        \# & Orbit Repre- &  Visualization  & $|\textrm{Aut}(W)|$  & Orbit \\
        &  sentative $W$ &   of $W$ &  & Length \\\toprule
         1 & $[1001:0001:0001]$ &\includegraphics[width=1.30cm]{Figures/counting/444/090101.png} & 2 & 768  \\\midrule
         \rowcolor{lime}2 & $[0101:0001:0110]$ &\includegraphics[width=1.30cm]{Figures/counting/444/050106.png} & 4 & 384  \\\midrule
         \rowcolor{lime}3 & $[0101:0100:0010]$ &\includegraphics[width=1.30cm]{Figures/counting/444/050402.png} & 4 & 384  \\\midrule
         \rowcolor{lime}4 & $[1011:1001:1111]$ &\includegraphics[width=1.30cm]{Figures/counting/444/110915.png} & 4 & 384  \\\midrule
         5 & $[1011:0001:0010]$ &\includegraphics[width=1.30cm]{Figures/counting/444/130102.png} & 6 & 256  \\\midrule
         6 & $[1101:1101:0100]$ &\includegraphics[width=1.30cm]{Figures/counting/444/131304.png} & 6 & 256  \\\midrule
         7 & $[1001:1100:0001]$ &\includegraphics[width=1.30cm]{Figures/counting/444/091201.png} & 8 & 192  \\\midrule
         \rowcolor{lime}8 & $[1110:0100:1111]$ &\includegraphics[width=1.30cm]{Figures/counting/444/140415.png} & 8 & 192  \\\midrule
         9 & $[0110:0111:0110]$ &\includegraphics[width=1.30cm]{Figures/counting/444/060706.png} & 8 & 192  \\\midrule
         \rowcolor{cyan}10 & $[0010:1010:0000]$ &\includegraphics[width=1.30cm]{Figures/counting/444/021000.png} & 8 & 192  \\\midrule
         \rowcolor{lime}11 & $[1000:0010:1111]$ &\includegraphics[width=1.30cm]{Figures/counting/444/080415.png} & 8 & 192  \\\midrule
         \rowcolor{cyan}12 & $[1010:0101:1011]$ &\includegraphics[width=1.30cm]{Figures/counting/444/100511.png} & 16 & 96  \\\midrule
         \rowcolor{cyan}13 & $[1100:1010:0000]$ &\includegraphics[width=1.30cm]{Figures/counting/444/121000.png} & 16 & 96  \\\bottomrule
         
    \end{tabular}
    \begin{tabular}{ccccc}
    \toprule
         \# & Orbit Repre- &  Visualization  & $|\textrm{Aut}(W)|$  & Orbit \\
         & sentative $W$ &   of $W$ &  & Length \\\toprule
         \rowcolor{cyan}14 & $[0101:1010:0110]$ &\includegraphics[width=1.30cm]{Figures/counting/444/051206.png} & 16 & 96  \\\midrule
         15 & $[0110:1001:0110]$ &\includegraphics[width=1.30cm]{Figures/counting/444/060906.png} & 24 & 64  \\\midrule
         \rowcolor{lime}16 & $[1111:0011:1100]$ &\includegraphics[width=1.30cm]{Figures/counting/444/150312.png} & 32 & 48  \\\midrule
         \rowcolor{cyan}17 & $[1100:0101:1010]$ &\includegraphics[width=1.30cm]{Figures/counting/444/120510.png} & 32 & 48  \\\midrule
         \rowcolor{lime}18 & $[0110:1001:0000]$ &\includegraphics[width=1.30cm]{Figures/counting/444/060900.png} & 32 & 48  \\\midrule
         \rowcolor{cyan}19 & $[1111:0011:0000]$ &\includegraphics[width=1.30cm]{Figures/counting/444/150300.png} & 32 & 48  \\\midrule
         \rowcolor{cyan}20 & $[1000:1111:1111]$ &\includegraphics[width=1.30cm]{Figures/counting/444/081515.png} & 32 & 48  \\\midrule
         \rowcolor{pink}21 & $[0000:0101:1010]$ &\includegraphics[width=1.30cm]{Figures/counting/444/000510.png} & 64 & 24  \\\midrule
         \rowcolor{pink}22 & $[1111:0101:0000]$ &\includegraphics[width=1.30cm]{Figures/counting/444/150500.png} & 64 & 24  \\\midrule
         \rowcolor{pink}23 & $[0101:1010:0101]$ &\includegraphics[width=1.30cm]{Figures/counting/444/051005.png} & 192 & 8  \\\midrule
         24 & $[0000:1111:1000]$ &\includegraphics[width=1.30cm]{Figures/counting/444/001508.png} & 192 & 8  \\\midrule
         \rowcolor{pink}25 & $[0000:0000:1111]$ &\includegraphics[width=1.30cm]{Figures/counting/444/000015.png} & 384 & 4  \\\midrule
         \rowcolor{pink}26 & $[1111:0000:1111]$ &\includegraphics[width=1.30cm]{Figures/counting/444/150015.png} & 384 & 4  \\\bottomrule
    \end{tabular}
    \label{tab:444_1}
\end{table}

\clearpage
\begin{table}
    \centering
    \caption{Table of $4\times4\times4$ Regular Plain cP-Weaves with rotational and translational symmetries. Red rows indicate that these cP-weaves can be obtained with repeating $2\times2\times2$ plain cP-weaves. Green rows indicate that these cP-weaves can be constructed from repeating $4\times4\times2$ plain cP-weaves. Blue rows indicate that these cP-weaves can be constructed from $4\times2\times2$ plain cP-weaves.}
    \begin{tabular}{ccccc}
    \toprule
        \# & Orbit Repre- &  Visualization  & $|\textrm{Aut}(W)|$  & Orbit \\
         & sentative $W$ &   of $W$ &  & Length \\\toprule
         1 & $[1001:0001:0001]$ &\includegraphics[width=1.5cm]{Figures/counting/444/090101.png} & 4 & 768  \\\midrule
         2 & $[1101:0001:0010]$ &\includegraphics[width=1.5cm]{Figures/counting/444/130102.png} & 6 & 512  \\\midrule
         \rowcolor{lime}3 & $[0101:0001:0110]$ &\includegraphics[width=1.5cm]{Figures/counting/444/050106.png} & 8 & 384  \\\midrule
         4 & $[1001:1100:0001]$ &\includegraphics[width=1.5cm]{Figures/counting/444/091201.png} & 8 & 384  \\\midrule
         \rowcolor{lime}5 & $[0101:0100:0010]$ &\includegraphics[width=1.5cm]{Figures/counting/444/050402.png} & 8 & 384  \\\midrule
         \rowcolor{lime}6 & $[0001:1111:1000]$ &\includegraphics[width=1.5cm]{Figures/counting/444/011508.png} & 8 & 384  \\\midrule
        \rowcolor{lime}7 & $[1011:1001:1111]$ &\includegraphics[width=1.5cm]{Figures/counting/444/110915.png} & 8 & 384  \\\midrule
         \rowcolor{cyan}8 & $[0010:1010:0000]$ &\includegraphics[width=1.5cm]{Figures/counting/444/021000.png} & 16 & 192 \\\midrule
         \rowcolor{cyan}9 & $[1010:0101:1011]$ &\includegraphics[width=1.5cm]{Figures/counting/444/100511.png} & 32 & 96  \\\midrule
         \rowcolor{cyan}10 & $[1100:1010:0000]$ &\includegraphics[width=1.5cm]{Figures/counting/444/121000.png} & 32 & 96  \\\bottomrule
     \end{tabular}
     \begin{tabular}{ccccc}
         \toprule
         \# & Orbit Repre- &  Visualization  & $|\textrm{Aut}(W)|$  & Orbit \\
         & sentative $W$ &   of $W$ &  & Length \\\toprule
         \rowcolor{lime}11 & $[1111:0011:1100]$ &\includegraphics[width=1.5cm]{Figures/counting/444/150312.png} & 32 & 96  \\\midrule
         \rowcolor{lime}12 & $[0101:1100:0110]$ &\includegraphics[width=1.5cm]{Figures/counting/444/051206.png} & 32 & 96  \\\midrule
         \rowcolor{cyan}13 & $[1111:0000:1101]$ &\includegraphics[width=1.5cm]{Figures/counting/444/150013.png} & 32 & 96 \\\midrule
         14 & $[0110:1001:0110]$ &\includegraphics[width=1.5cm]{Figures/counting/444/060906.png} & 48 & 64  \\\midrule
         \rowcolor{cyan}15 & $[1100:0101:1010]$ &\includegraphics[width=1.5cm]{Figures/counting/444/120510.png} & 64 & 48  \\\midrule
         \rowcolor{cyan}16 & $[1111:0011:0000]$ &\includegraphics[width=1.5cm]{Figures/counting/444/150300.png} & 64 & 48  \\\midrule
         \rowcolor{pink}17 & $[0000:0101:1010]$ &\includegraphics[width=1.5cm]{Figures/counting/444/000510.png} & 128 & 24  \\\midrule
         \rowcolor{pink}18 & $[1111:0101:0000]$ &\includegraphics[width=1.5cm]{Figures/counting/444/150500.png} & 128 & 24  \\\midrule
         \rowcolor{pink}19 & $[1111:0000:0000]$ &\includegraphics[width=1.5cm]{Figures/counting/444/150000.png} & 384 & 8  \\\midrule
         \rowcolor{pink}20 & $[0101:1010:0101]$ &\includegraphics[width=1.5cm]{Figures/counting/444/051005.png} & 384 & 8  \\\bottomrule
    \end{tabular}
    \label{tab:444_2}
\end{table}

\clearpage
\begin{table}
    \centering
    \caption{Number of Orbits for Regular Plain cP-Weaves with only rotational symmetries.}
    \begin{tabular}{cc}
    \midrule
        Weave Type & Number of Orbits \\\toprule
         $2\times2\times2$ & 10  \\\midrule
         $4\times4\times4$ & 240  \\\midrule
         $6\times6\times6$ & 11712  \\\midrule
         $8\times8\times8$ & 709376  \\\bottomrule
    \end{tabular}
    \label{tab:num1}
\end{table}
\clearpage

\begin{table}
    \centering
    \caption{Number of Orbits for Regular Plain cP-Weaves with rotational and mirror symmetries.}
    \begin{tabular}{cc}
    \midrule
        Weave Type & Number of Orbits \\\toprule
         $2\times2\times2$ & 6  \\\midrule
         $4\times4\times4$ & 124  \\\midrule
         $6\times6\times6$ & 5876  \\\midrule
         $8\times8\times8$ & 354808  \\\bottomrule
    \end{tabular}
    \label{tab:num2}
\end{table}
\clearpage
\begin{table}
    \centering
    \caption{Number of Orbits for Regular Plain cP-Weaves with rotational and translational symmetries.}
    \begin{tabular}{cc}
    \midrule
        Weave Type & Number of Orbits \\\toprule
         $2\times2\times2$ & 5  \\\midrule
         $4\times4\times4$ & 26  \\\midrule
         $6\times6\times6$ & 179  \\\midrule
         $8\times8\times8$ & 2504  \\\bottomrule
    \end{tabular}
    \label{tab:num3}
\end{table}
\clearpage
\begin{table}
    \centering
    \caption{Table of $4\times4\times4$ Regular Plain cP-Weaves with rotational, mirror and translational symmetries.}
    \begin{tabular}{cc}
    \midrule
        Weave Type & Number of Orbits \\\toprule
         $2\times2\times2$ & 4  \\\midrule
         $4\times4\times4$ & 20  \\\midrule
         $6\times6\times6$ & 129  \\\midrule
         $8\times8\times8$ & 1468  \\\bottomrule
    \end{tabular}
    \label{tab:num4}
\end{table}


\end{document}
